% Fichier généré automatiquement par src/analysis : ne pas modifier à la main.
\begin{table}[H]
\centering\small
\begin{tabular}{@{}lccc@{}}
\toprule
\textbf{Modèle} & \textbf{Exactitude [IC 95\,\%]} & \textbf{F1 [IC 95\,\%]} & \textbf{AUC [IC 95\,\%]} \\
\midrule
Baseline CNN & 0{,}836 [0{,}824 ; 0{,}847] & 0{,}837 [0{,}824 ; 0{,}848] & 0{,}921 [0{,}912 ; 0{,}929] \\
EfficientNet-V2-S & 0{,}971 [0{,}966 ; 0{,}977] & 0{,}971 [0{,}966 ; 0{,}977] & 0{,}996 [0{,}995 ; 0{,}997] \\
ViT-B/16 & 0{,}997 [0{,}996 ; 0{,}999] & 0{,}997 [0{,}996 ; 0{,}999] & 1{,}000 [1{,}000 ; 1{,}000] \\
ViT-B/16 (\texttt{scale}) & 0{,}999 [0{,}997 ; 1{,}000] & 0{,}999 [0{,}997 ; 1{,}000] & 1{,}000 [1{,}000 ; 1{,}000] \\
\bottomrule
\end{tabular}
\caption{Métriques sur le jeu de test avec intervalles de confiance à 95\,\% (bootstrap percentile, 2000 ré-échantillonnages des 4\,000 images, graine d'entraînement 42).}
\end{table}
\begin{table}[H]
\centering\small
\begin{tabular}{@{}lccc@{}}
\toprule
\textbf{Comparaison (A contre B)} & \textbf{Seul A a raison} & \textbf{Seul B a raison} & \textbf{p (McNemar exact)} \\
\midrule
ViT-B/16 (\texttt{scale}) contre Baseline CNN & 654 & 4 & $<10^{-4}$ \\
ViT-B/16 (\texttt{scale}) contre EfficientNet-V2-S & 114 & 5 & $<10^{-4}$ \\
ViT-B/16 (\texttt{scale}) contre ViT-B/16 & 8 & 3 & 0{,}2266 \\
\bottomrule
\end{tabular}
\caption{Test de McNemar sur les mêmes images de test : l'écart entre deux modèles est significatif si $p<0{,}05$.}
\end{table}
L'avantage de ViT-B/16 (\texttt{scale}) sur Baseline CNN, EfficientNet-V2-S est statistiquement significatif (McNemar, $p<0{,}05$) : il ne s'explique pas par le hasard de l'échantillon de test. En revanche, l'écart avec ViT-B/16 n'est pas significatif : les deux modèles sont équivalents sur ce jeu de test. 

\begin{table}[H]
\centering\small
\begin{tabular}{@{}lcccc@{}}
\toprule
\textbf{Modèle} & \textbf{Graines} & \textbf{Exactitude} & \textbf{F1} & \textbf{AUC} \\
\midrule
EfficientNet-V2-S & 1, 2, 42 & 0{,}973 $\pm$ 0{,}001 & 0{,}973 $\pm$ 0{,}001 & 0{,}997 $\pm$ 0{,}001 \\
ViT-B/16 & 1, 2, 42 & 0{,}996 $\pm$ 0{,}001 & 0{,}997 $\pm$ 0{,}001 & 1{,}000 $\pm$ 0{,}000 \\
\bottomrule
\end{tabular}
\caption{Variabilité entre entraînements : moyenne $\pm$ écart-type sur le jeu de test pour plusieurs graines.}
\end{table}
